\section{What if families look beyond the shortlist?}\label{sec:response}
The platform shows five discretionary courses. Families can ignore it, restore an omitted choice, or look for a different school. We can observe what families sought after the actual regular round and how a risk warning changed applications in an earlier field study. We use each source for the question it can answer.

\paragraph{Adding a school before the regular round clears}
Near a 30\% predicted nonplacement threshold, a Chilean warning raised the probability of adding a school from about 19.9\% to 41.5\%. Most of the increase involved appending a school \citep{smartplatforms,warningdata}. We use each probability as a scenario for pupils who have an undisplayed school on their submitted list. In each scenario the same pupils are willing to respond under both menus. We select their highest-ranked course at a school entirely absent from that menu and run the complete institutional allocation again. In Table~\ref{tab:behavior-main}, coverage now credits every school pursued. The original displays still differ by 4,638 hits throughout.

\begin{table}[ht]\centering\small
\caption{What happens when families pursue more than the five recommendations? Informed minus blind, among the same 1,969,120 register-linked pupils. An April hit means the historically attended school is in the pursued set. Fractional counts are means over eight paired repetitions; other rows are deterministic.}\label{tab:behavior-main}
\begin{tabular}{@{}lrr@{}}\toprule
How families pursue schools & April hits & Linked placements\\\midrule
Use the displayed menu & $+4,638$ & $-1,670$\\
Append first omitted school, 19.9\% participation & $+3,841.5$ & $-1,664.3$\\
Append first omitted school, 41.5\% participation & $+2,971.4$ & $-1,656.0$\\
Restore that school in submitted order, 41.5\% & $+2,971.4$ & $-1,340.8$\\
Append courses the same pupil explicitly listed again & $+4,733$ & $-1,653$\\
Append all omitted courses & $0$ & $-1,234$\\
Restore all eligible courses in submitted order & $0$ & $0$\\\bottomrule
\end{tabular}
\end{table}

\paragraph{Who chooses to keep looking?}
Giving everyone the same response probability can conceal a reversal. Let the same families respond under both policies, and fix their number at 19.9\% of pupils able to add a school under either. If they append their first omitted school, the sharp range of the April-hit difference is $[-179,5{,}412]$. At the lower endpoint the rank-informed rule both misses 179 more eventual schools and places 1,669 fewer linked pupils. If instead the added course has the lowest demand relative to vacancies among omitted schools, even the least favorable common set of responders leaves a gain of 920 hits. The placement gap for that set is $-1{,}612$. The arithmetic behind the bounds is simple: a respondent matters for the comparison only when adding the course turns a miss into a hit under exactly one policy. The placement counts require national re-clearing. Appendix~\ref{sec:behavior-tests} gives the proof, the constructed responder sets and the sensitivity to unequal uptake across policies.

\paragraph{Where families actually went}
Among 154,988 same-grade reapplicants, 139,190 named a school absent from their entire regular report. Only 20,698 had an eligible original course omitted by either hypothetical menu; 1,380 explicitly listed such a course again. Of 10,296 with an omitted course showing a positive complementary vacancy, 1,368 reapplied to one. Appending those exact repeat choices gives 4,733 more hits and 1,653 fewer linked placements. A different check uses the identities of the 154,988 actual returners but lets each append the highest-ranked school missing from their menu. Only 20,501 have such a school under either policy. The comparison gives 4,572 more hits and 1,665 fewer linked placements. The historical record points to a larger margin: families went looking for schools they had never listed.

To test that margin, we keep every simulated first-round place and give unplaced pupils access to the courses at new schools they actually named in the complementary round. Both arms face the same published second-round vacancies. Ignoring priorities, we find the largest possible number of extra placements in each arm. Across all applicants the two menus can add 63,571 and 63,847 recipients respectively; this shrinks the first-stage gap from 1,834 to 1,558. Giving register-linked pupils first claim on these vacancies leaves 1,419 fewer linked placements under the informed rule and 4,752 more pursued sets containing the April school. Families might reapply differently if their initial offers changed. These numbers describe how far their \emph{recorded} search could reach in a common inventory. They do not predict the offers of an actual second round.

In the rental simulations, appending every omitted feasible apartment brings the enriched-market placement difference to $+0.28$ percentage points (paired block interval $[-0.07,0.65]$). That interval includes zero; the corresponding ordinary-market interval remains below zero. The full set of comparisons appears in Appendix~\ref{sec:behavior-tests}.
